\par\noindent\textbf{(1)} 电流为 \(I_i=\dot q_i\)，磁能 \(T_L=\tfrac12(L_1I_1^2+L_2I_2^2+2MI_1I_2)\)，电场储能 \(U_C=q_1^2/(2C_1)+q_2^2/(2C_2)\)。Lagrangian取 \(L=T_L-U_C\)。Euler方程为 \(L_1\ddot q_1+M\ddot q_2+q_1/C_1=0\) 与 \(M\ddot q_1+L_2\ddot q_2+q_2/C_2=0\)。
\par\noindent\textbf{(2)} 取 \(q_j=Q_je^{-i\omega t}\)，非零模态要求 \(\det\begin{pmatrix}C_1^{-1}-\omega^2L_1&-\omega^2M\\-\omega^2M&C_2^{-1}-\omega^2L_2\end{pmatrix}=0\)。展开得到 \((L_1L_2-M^2)\omega^4-(L_1/C_2+L_2/C_1)\omega^2+1/(C_1C_2)=0\)。
\par\noindent\textbf{(3)} 对 \(L_1=L_2=L\)、\(C_1=C_2=C\)，同相 \(Q_1=Q_2\) 满足 \(\omega_+^2=1/[(L+M)C]\)，反相 \(Q_1=-Q_2\) 满足 \(\omega_-^2=1/[(L-M)C]\)。正互感增加同相磁能，因而同相模频率降低，反相模频率升高。
\par\noindent\textbf{(4)} 储磁能矩阵 \(\begin{pmatrix}L_1&M\\M&L_2\end{pmatrix}\) 必须正定，即 \(L_1L_2>M^2\)。在 \(|M|\to\sqrt{L_1L_2}\) 极限，一个组合的磁惯量趋零、相应理想模频率发散；这表明模型的零漏感奇异性，而非可实现无限频率。
